\subsection*{Operaciones 1 Base 2}
  
Se define un marco de referencia y a el se fija la siguiente base $\{\Vec{e}_1,\Vec{e}_2,\Vec{e}_3\}$ que es orto-normal, resuelva las siguientes operaciones:
\begin{multicols}{2}
\begin{enumerate}
    \item$\Vec{e}_1\cdot\Vec{e}_2=$
    %\item$a\Vec{e}_1\cdot \Vec{e}_3=$
    \item$\Vec{e}_2\cdot a\Vec{e}_2=$
    \item$\Vec{e}_3\times \Vec{e}_1=$
    %\item$\|a\Vec{e}_3\times b\Vec{e}_3\|=$
    \item$\|a\Vec{e}_1\times b\Vec{e}_2\|=$
    %\item$\Vec{e}_1\cdot (\Vec{e}_2+a\Vec{e}_3)=$
    \item$\Vec{e}_1\cdot (\Vec{e}_2+a\Vec{e}_1 )=$
    %\item$(a\Vec{e}_1+d\Vec{e}_3)\cdot (b\Vec{e}_2+c\Vec{e}_3)=$
    \item$\|(a\Vec{e}_1+d\Vec{e}_3)\times (b\Vec{e}_2+c\Vec{e}_3)\|=$
\end{enumerate}
\end{multicols}
\subsection*{Operaciones 2 Bases}

Se tienen 2 marcos de referencia a cada uno  se le fija una base orto-normal, al primero cartesiana $\{\Vec{e}_x,\Vec{e}_y,\Vec{e}_z\}$ y al segundo una cilíndrica $\{\Vec{e}_r,\Vec{e}_\theta,\Vec{e}_z\}$, se conoce la transformación de una base a otra: 
\begin{equation*}
\begin{split}
    \Vec{e}_r&=\cos(\theta)\Vec{e}_x+\sin(\theta)\Vec{e}_y\\
    \Vec{e}_\theta&=-\sin(\theta)\Vec{e}_x+\cos(\theta)\Vec{e}_y\\
    \Vec{e}_z&=\Vec{e}_z
\end{split}
\end{equation*}

\begin{multicols}{2}
\begin{enumerate}
    \item exprese el vector $\Vec{e}_z$ en la base cilíndrica 
    \item exprese el vector $\Vec{e}_x$ en la base cilíndrica 
    \item exprese el vector $\Vec{e}_y$ en la base cilíndrica 
    \item $\Vec{e}_z\times\Vec{e}_\theta=$
   % \item $a\Vec{e}_z\times b\Vec{e}_r=$
    \item $a\Vec{e}_r\cdot b\Vec{e}_x=$
    \item $l\Vec{e}_r\times(-mg\Vec{e}_y)=$
    \item $\Vec{e}_z\cdot\Vec{e}_\theta=$
    \item $\|a\Vec{e}_y\times\Vec{e}_\theta+\Vec{e}_x\|=$
   % \item $a\Vec{e}_x\times\Vec{e}_r=$
   % \item $l\Vec{e}_r\times-mg\Vec{e}_y=$
    
\end{enumerate}
\end{multicols}


